\section{Definitions and Setup for the Iterative Construction}\label{sec:input-statement}

An arrangement of lines is called \emph{simple} if no point belongs to more
than two of its lines. We work with simple affine arrangements in
$\mathbb{R}^2$ and write $a_3(\cA)$ for the number of bounded triangular
faces of an arrangement~$\cA$. A distinguished line $Y_0$ \emph{touches} a
triangle if it contains one of its sides.

Throughout the paper we use the best affine upper bounds for odd $n$
from \cite[Theorem~1]{blanc2011}: for every simple affine arrangement
of $n$ lines or pseudolines,
\begin{equation}\label{eq:upper-bound}
  a_3(\cA)\le
  \begin{cases}
    (n(n-2)-2)/3 & \text{if } n\equiv 1\pmod{6}, \\
    n(n-2)/3 & \text{if } n\equiv 3,5\pmod{6}.
  \end{cases}
\end{equation}

We use the following affine formulation of the iterative step from
\cite[Proposition~3.1 and Remark~3.2]{bartholdi2007}.

\begin{proposition}[Affine iterative step]\label{prop:iterative-step}
Let $n\ge 2$ be an even number, and let
\[
  \cA=\{Y_0,L_1,\dots,L_n\}
\]
be a simple affine arrangement of $n+1$ lines, given by the equations
\[
  Y_0:\ y=0,
  \qquad
  L_i:\ y=m_i(x-a_i).
\]
Assume that
\[
  \{a_1,\dots,a_{n-2}\}=\{\pm\tan(k\pi/n)\mid k=1,\dots,n/2-1\},
\]
and
\[
  -\frac{1}{n}<a_{n-1}<0<a_n<\frac{1}{n},
\]
and that $Y_0$ touches exactly $n-1$ triangles of the arrangement $\cA$.
Then there exist $n$ lines $M_1,\dots,M_n$ such that
the affine arrangement
\[
  \cB=\{Y_0,L_1,\dots,L_n,M_1,\dots,M_n\}
\]
of $2n+1$ lines is simple and has exactly $n^2$ more
triangles than $\cA$, and $Y_0$ touches exactly $2n-1$ triangles of $\cB$.
\end{proposition}

\begin{remark}[Iteration condition]\label{rem:iteration-condition}
If moreover $|a_{n-1}|$ and $|a_n|$ are smaller than $1/(2n)$, then the new
arrangement $\cB$ again satisfies the hypotheses of the proposition. Hence the
construction can be iterated provided these two special parameters are chosen
arbitrarily small.
\end{remark}

For the base configuration studied in this paper
(see Figure~\ref{fig:base-configuration}), the relevant value is $n=18$,
so it has $19$ lines in total. In Section~\ref{sec:main-result},
we prove the existence of a straight-line affine arrangement
$
  \{Y_0,L_1,\dots,L_{18}\}
$
with
\[
  \{a_1,\dots,a_{18}\}=\{\pm\tan(k\pi/18)\mid k=1,\dots,8\} \cup \{-\eps,+\eps\},
\]
where $\eps>0$ is small enough and the hypotheses of the proposition and
remark above hold.
